\section{Experimental Setup}
\label{sec:experiments}

We design five experiments to test our contamination-correction hypothesis.

\subsection{Research Questions}

\textbf{RQ1 (Existence):} Does deduplication produce a statistically significant per-benchmark
accuracy differential at token-count-matched checkpoints?

\textbf{RQ2 (Correlation):} Does the per-benchmark accuracy differential correlate positively
with estimated n-gram contamination?

\textbf{RQ3 (Mechanism --- min-$k$\%):} Do Pile-trained models show higher min-$k$\% scores,
consistent with near-memorization?

\textbf{RQ4 (Methodology):} Does token-count matching recover a stronger contamination signal
than step-matching?

\textbf{RQ5 (Documents):} Do removed documents show higher n-gram overlap with benchmarks
than retained documents?

\subsection{Evaluation Settings}

\begin{table}[h]
\centering
\caption{Benchmark evaluation settings and estimated contamination rates.}
\label{tab:eval_settings}
\begin{tabular}{llcc}
\toprule
Benchmark & Task Type & Shots & 13-gram Contam. (Pile) \\
\midrule
MMLU & General knowledge & 5 & 5.5\% \\
HellaSwag & Commonsense & 0 & 20.0\% \\
ARC-Challenge & Science & 25 & 8.5\% \\
WinoGrande & Commonsense & 5 & 2.5\% \\
\bottomrule
\end{tabular}
\end{table}

\textbf{Statistical tests:} The existence analysis (RQ1) uses paired $t$-test across 4 model sizes, Bonferroni
$\alpha = 0.0125$. The correlation analysis (RQ2) uses Pearson $r$ at $\alpha = 0.05$ with bootstrap CI.
The memorization analysis (RQ3) uses paired $t$-test (one-tailed, Pile $>$ dedup). The document overlap analysis (RQ5) uses Mann-Whitney $U$.

All evaluations use lm-evaluation-harness greedy decoding, half-precision (float16) GPU inference.
